% LỜI TỰA --- Quyển XI
\chapter*{Lời Tựa}
\addcontentsline{toc}{chapter}{Lời Tựa}
\markboth{Lời Tựa}{Lời Tựa}

\begin{truyen}{Lời Tựa}{Traps Young People Set for Themselves}
\chuhoa{G}{iáo viên bắt đầu bộ sách này với một câu hỏi đơn giản: làm thế nào để mỗi buổi học có thể để lại thứ gì đó đáng nhớ hơn là bài thi?}
\textit{(A teacher began this series with a simple question: how can each class leave behind something more memorable than a test?)}

Hiểu tên của cái bẫy là cách đầu tiên để thoát khỏi nó.
\textit{(Knowing the name of the trap is the first step to escaping it.)}

Quyển mười một đặt tên cho những cái bẫy mà người trẻ thường tự giăng: tự mãn, cái tôi quá lớn, so sánh thường xuyên, trì hoãn mãn tính. Không ai miễn nhiễm~--- nhưng ai đọc quyển này với tâm thế thành thật sẽ thấy mình trong ít nhất một câu chuyện.
\textit{(Volume eleven names the traps young people commonly set for themselves: complacency, oversized ego, constant comparison, chronic procrastination. No one is immune~--- but anyone who reads this volume with honesty will find themselves in at least one story.)}

Bộ sách <<Truyện Tích Cảm Hứng Song Ngữ>> gồm 24 quyển, mỗi quyển một chủ đề riêng. Quyển XI này mang chủ đề: \textbf{tự mãn, cái tôi, so sánh, trì hoãn}.
\textit{(The series ``Inspiring Bilingual Tales'' contains 24 volumes, each with its own theme. Volume XI focuses on: \textbf{tự mãn, cái tôi, so sánh, trì hoãn}.})
\end{truyen}

\begin{baihoc}
Bẫy tự bước vào đau hơn bẫy người khác đặt~--- vì nó thường đi kèm với sự xấu hổ.
\textit{(Traps we walk into ourselves hurt more than those others set~--- because they usually come with shame.)}
\end{baihoc}

\begin{ghinhoanh}
\textbf{Short English to remember:}

Name your trap: that is the first step out.

Self-awareness is the sharpest tool against self-sabotage.
\end{ghinhoanh}

\begin{tuhocnhanh}
1. Em đang bước vào cái bẫy nào mà vẫn chưa chịu thừa nhận? \textit{(What trap are you walking into that you haven't admitted yet?)}

2. Điều gì khiến em trì hoãn nhiều nhất? Em đã thử cách nào để thoát chưa? \textit{(What do you procrastinate most on? Have you tried any way out?)}

3. Trong những cái bẫy được kể, cái nào em thấy phổ biến nhất ở bạn bè? \textit{(Among the traps described, which do you see most commonly in your friends?)}
\end{tuhocnhanh}

\begin{center}
\textit{\color{orange}``Thoát bẫy không phải là chuyện xấu hổ~--- đó là chuyện dũng cảm.''}
\end{center}

\begin{flushright}
\textbf{Tôi Nguyễn Văn Sang}\\
\textit{GV THPT Nguyễn Hữu Cảnh - TP HCM}
\end{flushright}

\ngancach
